% 見本の章（修正版）．忠実版を写して直したもの．直した箇所の先頭に \fixat{ID} を置く．
\chapter{距離}\label{ch:1}
\origpage{1}

\section{定義}\label{sec:1.1}

\begin{definition}\label{thm:1.1.1}
集合 $X$ 上の関数 $d\colon X\times X\to[0,\infty)$ が，任意の $x,y,z\in X$ に対して次の 3 つをみたすとき，$d$ を $X$ 上の\term{距離}という．
\begin{equation}\label{eq:1.1}
\begin{gathered}
d(x,y)=0 \iff x=y,\qquad d(x,y)=d(y,x),\\
d(x,z)\le d(x,y)\fixat{1-01}+d(y,z).
\end{gathered}
\end{equation}
\end{definition}

\begin{proposition}\label{thm:1.1.2}
任意の $x,y,z\in X$ に対して $|d(x,z)-d(y,z)|\le d(x,y)$ が成り立つ．
\end{proposition}
\fixat{1-02}\begin{proof}\supplied{%
三角不等式 \beqref{eq:1.1} より $d(x,z)\le d(x,y)+d(y,z)$ なので，$d(x,z)-d(y,z)\le d(x,y)$ である．
$x$ と $y$ を入れ替えれば $d(y,z)-d(x,z)\le d(y,x)$ である．
対称性より $d(y,x)=d(x,y)$ なので，2 つを合わせて主張を得る．}
\end{proof}

\begin{proposition}\label{thm:1.1.3}
$d$ が $X$ 上の距離なら，$d'=\min\{d,1\}$ も $X$ 上の距離である．
\end{proposition}
\fixat{1-03}\begin{proof}\rewritten{%
三角不等式だけを示す．関数 $t\mapsto\min\{t,1\}$ は単調で劣加法的である．
よって $d'(x,z)\le\min\{d(x,y)+d(y,z),1\}\le d'(x,y)+d'(y,z)$ となる．
\origpage{2}%
ここで右辺の各項は $d'$ で置き換えてもよいことに注意する．}
\end{proof}
\fixskip{equation}{1}% 書き換えで消した (1.2) の分

後の式の番号は原書と同じである:
\begin{equation}\label{eq:1.3}
d'(x,y)\le d(x,y).
\end{equation}

\begin{example}\label{thm:1.1.4}
$\mathbb{R}$ 上で $d(x,y)=|x-y|$ は距離である\fixat{1-05}\supplied{（三角不等式は絶対値の三角不等式）}．
\end{example}

\fixat{1-09}\supplied{%
\begin{addproposition}\label{rev:1.1.4a}
任意の $x,y,z,w\in X$ に対して
\begin{equation*}\addtag\label{rev:1.3a}
|d(x,y)-d(z,w)|\le d(x,z)+d(y,w).
\end{equation*}
\end{addproposition}
\begin{proof}
命題 \bref{thm:1.1.2} を 2 回使えば \beqref{rev:1.3a} を得る．
\end{proof}}

\section{\texorpdfstring{$\varepsilon$}{ε} 近傍}\label{sec:1.2}

\begin{theorem}\label{thm:1.2.1}
$x\in X$ と $\varepsilon>0$ に対し，集合 $B(x,\varepsilon)=\{y\in X\mid d(x,y)<\varepsilon\}$ は開集合である\fixat{1-06}\supplied{（注：開集合の定義は本書にない．ここでは通常の定義を仮定する）}．
\end{theorem}

次の章の定理 \bref{thm:2.1.1} で，命題 \bref{thm:1.1.2} を使う．\fixat{1-08}印の付いた文．
